\subsection{Acceso, actualización y ejecución de compuertas}

La aplicación de compuertas sobre el backend comprimido se organizó de acuerdo con el
patrón de acceso que cada operación induce sobre las amplitudes. En términos
generales, la implementación distingue entre operaciones cuyos pares afectados pueden
resolverse dentro de un mismo bloque y operaciones que involucran amplitudes
distribuidas en bloques distintos. Esta distinción permite adaptar la ruta de ejecución al
tipo de acceso requerido, en lugar de tratar todas las compuertas como si implicaran la
materialización completa del vector.

Cuando la operación permanece dentro de un bloque, el backend adquiere ese bloque,
opera sobre su buffer descomprimido y marca el contenido como modificado cuando
corresponde. Cuando la compuerta involucra amplitudes ubicadas en bloques distintos,
la ejecución requiere adquirir simultáneamente los bloques implicados y coordinar la
actualización sobre ambos buffers. Esta lógica se articula a través del propio backend y
de componentes auxiliares como \texttt{PairKernelExecutor} y
\texttt{GateApplyEngine}, que separan rutas optimizadas para pares de amplitudes de una
ruta más general para compuertas de aplicación genérica.

La implementación incluye además un mecanismo de agrupación operativa para delimitar
el inicio y el cierre de lotes de trabajo. Este permite agrupar operaciones consecutivas y
diferir ciertas escrituras mientras los bloques implicados permanecen activos. Así se
reducen recomposiciones innecesarias y se mejora la interacción entre acceso al backend,
reutilización de buffers y persistencia diferida de datos.

Desde el punto de vista funcional, esta organización hace posible que la lógica del
simulador siga operando sobre amplitudes complejas y compuertas cuánticas del mismo
modo que en una versión no comprimida, mientras la capa de almacenamiento resuelve
de forma interna la recuperación, actualización y escritura de los bloques necesarios para
cada paso de la simulación.
